\chapter{Conclusiones}

Este capítulo recoge la valoración global del trabajo realizado, analizando el grado
de cumplimiento de los objetivos planteados al inicio del proyecto y reflexionando
sobre las decisiones técnicas adoptadas a lo largo del desarrollo. Se incluye además
una declaración sobre el uso de herramientas de inteligencia artificial generativa
durante el proceso, y se proponen las líneas de trabajo futuro más relevantes para
continuar evolucionando el proyecto más allá del alcance definido en este TFG.

El desarrollo de \textit{NeuroForge} ha permitido alcanzar de forma satisfactoria los
objetivos principales planteados al inicio del proyecto. En primer lugar, se ha
diseñado e implementado con éxito la lógica fundamental del videojuego, consolidando
una estructura de tablero funcional, la definición precisa de las piezas y las reglas
que rigen su movimiento y combate. Sobre esta base, se ha desarrollado un robusto
sistema de control de turnos y gestión del estado de la partida, lo que garantiza la
ejecución de enfrentamientos completos bajo el estricto cumplimiento de las mecánicas
diseñadas. Todo este proceso ha propiciado una valiosa adquisición de experiencia
práctica en el uso del motor de videojuegos Godot y en el desarrollo de sistemas
interactivos empleando C\# como lenguaje principal. Como resultado, la experiencia de
juego es coherente, estable y accesible. La interfaz de usuario alcanza el nivel de
pulido propuesto, todas las pantallas están implementadas, el flujo de navegación es
intuitivo y el apartado audiovisual refuerza adecuadamente la ambientación de ciencia
ficción del proyecto. La arquitectura en tres capas adoptada ha demostrado ser una
decisión acertada, manteniendo la lógica del videojuego completamente desacoplada del
motor y facilitando tanto las pruebas como el mantenimiento del código a lo largo de
las iteraciones. El desarrollo del bot inteligente mediante aprendizaje por refuerzo
constituyó el objetivo técnicamente más ambicioso del proyecto y ha sido completado
satisfactoriamente. El entorno de entrenamiento basado en Gymnasium replica con
fidelidad las reglas del videojuego, incluyendo la gestión de información oculta y
las reglas especiales de combate, y el agente entrenado mediante MaskablePPO y
estrategia de \textit{self-play} ha demostrado un comportamiento estratégico
coherente, siendo capaz de gestionar la incertidumbre inherente al juego sin necesidad
de una función de evaluación diseñada manualmente. Su integración con Godot a través
de un subproceso Python desacoplado del motor ha resultado en una solución técnicamente
sólida y fácilmente actualizable.

Durante el desarrollo de este proyecto se ha hecho uso de herramientas de inteligencia
artificial generativa como apoyo en determinadas tareas, sin que en ningún caso hayan
sustituido la toma de decisiones técnicas ni el trabajo de diseño e implementación
propios. En el ámbito de la documentación, la IA se empleó principalmente como
asistente para la redacción en \LaTeX{}, para resolver dudas sobre sintaxis específica,
formatear correctamente tablas o figuras, y transformar texto redactado a mano en
código \LaTeX{} bien estructurado. Esto permitió agilizar la fase de escritura de la
memoria sin delegar en ningún momento el contenido ni el criterio técnico de lo
escrito. En el ámbito del código, su uso se limitó a tareas de soporte, refactorización
de fragmentos para mejorar su legibilidad, consultas puntuales sobre la API de C\# o
sobre patrones de uso específicos de Godot, y en algún caso la sugerencia de
estructuras alternativas para organizar mejor un bloque de lógica ya diseñado. La
arquitectura del sistema, las decisiones de diseño y la implementación de las mecánicas
del videojuego fueron elaboradas de forma íntegramente propia.

El estado actual del proyecto constituye una base sólida sobre la que es posible crecer
en múltiples direcciones. Una de las fortalezas de la arquitectura implementada es que
el tablero no está codificado de forma rígida en la lógica del videojuego, sino que se
define como una matriz que Godot interpreta directamente desde el editor, lo que
significa que crear una variante con dimensiones distintas, una distribución diferente
de zonas intransitables o un conjunto de piezas modificado requiere únicamente editar
esa configuración sin tocar el código, abriendo la puerta a modos de juego alternativos
con un coste de desarrollo muy reducido. Con un esfuerzo algo mayor, la estructura de
datos del estado de la partida se presta de forma natural a la serialización, por lo
que implementar un sistema de guardado y carga consistiría esencialmente en serializar
ese estado a un archivo JSON y deserializarlo al reanudar sin necesidad de modificar la
arquitectura actual. En esa misma línea, con el agente inteligente ya entrenado e
integrado, resultaría relativamente sencillo exponer distintos niveles de dificultad
mediante modelos entrenados durante diferente número de iteraciones, ajustando el grado
de aleatoriedad de la política o combinando el bot inteligente con el bot aleatorio en
proporciones variables. Con las bases del sistema implementadas, añadir nuevo contenido
tiene también un coste relativamente bajo: modificadores opcionales como alterar la
regla de empate o introducir un temporizador por turno, nuevos tipos de piezas como una
torre de radar que permita revelar una pieza enemiga tras cierto número de turnos, o
nuevos tipos de casillas como teletransportadores o zonas de peligro periódico,
enriquecerían la experiencia sin necesidad de rediseñar el sistema de base. Finalmente,
respecto a la publicación, antes de plantear una distribución comercial resultaría
conveniente publicar el videojuego en \textit{itch.io} para obtener retroalimentación
real de jugadores sin coste de entrada. Una vez incorporadas las mejoras anteriores,
especialmente el bot con niveles de dificultad, el sistema de guardado y una variedad
de tableros adicionales, el producto estaría en condiciones de justificar su presencia
en \textit{Steam} al precio de 1,99~€ analizado en el capítulo de planificación, siendo
la sustitución de los recursos visuales actuales por arte original el paso que
marcaría la diferencia entre un prototipo académico y un producto independiente con
identidad propia.